%%%%%%%%%%%%% Common preamble %%%%%%%%%%%%
\documentclass[
    12pt
    ,a4paper
    ,oneside
%    ,chapterprefix
    ,russian
]{memoir}
%\usepackage{subfiles}
\usepackage{amssymb} % Essentials to be included before xunicode
\usepackage{ifxetex}
%%%%%%%%%%%%%%%%%%%%%%%%%%%%%%%%%%%%%%%%%%

\ifxetex

\usepackage{fontspec}
\usepackage{xunicode}
\usepackage{xltxtra}
\usepackage{polyglossia}
\setromanfont[Mapping=tex-text]{Liberation Serif}
\setsansfont[Mapping=tex-text]{Liberation Sans}
\setmonofont[Mapping=tex-text]{Liberation Mono}
\setdefaultlanguage{russian}

\else % not XeTeX

\end{document}

\fi % not XeTex

%%%%%%%%%%%%% Common preamble %%%%%%%%%%%%

%\usepackage{graphicx}
\usepackage{textcomp}

\usepackage{indentfirst}
\frenchspacing
\clubpenalty=10000
\widowpenalty=10000

\sloppy

\newcommand{\mybreakfleurons}{
    {\textasteriskcentered}
    \hspace{1em}
    {\textasteriskcentered}
    \hspace{1em}
    {\textasteriskcentered}
}

\newcommand{\myplainbreak}{
    \medskip
}

\newcommand{\myfancybreak}{
    \fancybreak*{\mybreakfleurons}
    %\centering{\bigskip{\mybreakfleurons}} % scrbook mode
}

\newcommand{\mydatestamp}{
    \bigskip   % Temporary solution, until paragraph spacing is discovered.
    \emph
}

\newcommand{\mds}{\mydatestamp}

%%%%%%%%%%%%%%%%%%%%%%%%%%%%%%%%%%%%%%%%%%

\begin{document}
    \begin{titlingpage}
        \title{Веркольский дневник}
        \author{Алексей Орлов}
        \date{~}
        \maketitle
    \end{titlingpage}
    
    \chapter*{2017}
    \section*{январь}
    
    \mds{2017-01-27, пятница}. В поезде Санкт-Петербург-Архангельск душно. Вентиляция исправна~---~просто ночью ее включают на десять минут раз в три часа. На случайность не похоже: в вагоне полно народу в футболках и шортиках. Зимой. Опыт, не иначе. Любопытно, что после условного <<подъема>> климат становится нормальным. Вот когда бы спать, но уже все. В поезде Архангельск-Карпогоры похожая история, но тут можно поверить, что вентиляция и впрямь не работает. Вагон старого образца, с форточками, но все заперто на ключ.
    
    Надо сказать, в районе вокзала Архангельск совершенно неинтересен, чтоб не сказать гнусен. Пересидевший промежуток между поездами в зале ожидания не теряет ничего.
    
    Маршрутка от Карпогор до деревни Веркола идет довольно долго, и все под гору. Не ожидал такого интересного рельефа.
    
    Последние два километра пешком от Верколы до монастыря через Пинегу возмещают все потери: над головой звездное небо, куда интереснее того, каким пользуется Пулковская обсерватория, в полной темноте фонарик работает не хуже автомобильных фар. Если не оборачиваться, вполне сойдет за снежную пустыню. Впечатление волшебное, притом еще не слишком холодно. Реку подо льдом опознать можно только по тонким деревянным вешкам с катафотами, каких на грунте никто ставить не будет. По льду идет нормальная дорога, со следами не только протекторов, но и гусениц. Крепкий лед, хотя температура, вроде бы, не опускается ниже минус пятнадцати, и то в отдельные дни.
    
    \mds{2017-01-28, суббота}. Первый день в монастыре. Прошедшая ночь не в счет: меня просто уложили спать в келье на пять человек. Сейчас нас трое: я, Роман Киселев и Сергей Овчинников. На новом месте мелкие неурядицы: монастырский этикет не похож на офисный. Двери келий не запираются, надо стучаться с молитвой. Если не слышно ответа <<аминь>>, келья считается запертой, а хозяин~---~отсутствующим. Из знакомого: мобильные устройства и ноутбуки запрещены, но у всех есть, лежат и используются открыто.
    
    Набрал сдуру в каптерке всякой одежды для грязной работы, а складывать негде, теперь надо возвращать почти все, и сумку собственных вещей сдать туда же. Места для личных вещей в келье мало.
    
    Работал немного на дровах, и таскал мешки с песком. Все, что можно сказать о физической форме: нести мешок могу, а поднимать его мне на плечи приходится другим.
    
    Еда неплохая, и волу молотящему рта никто не заграждает.
    
    Котов пока два, Моня черной масти, и Рыжий. Собак, вроде бы, нет ни одной. Ночью, когда я пришел, никто не лаял.
    
    \mds{2017-01-29, воскресенье}. Познакомился с библиотекарем о. Матфеем. Поговорил, как половчее взяться за изучение церковнославянского языка, мало что узнал нового, зато получил Библию на церковнославянском и книгу святителя Игнатия Брянчанинова <<Аскетические опыты>>, с напутствием, что книга эта является лучшим введением в православие.
    
    Обнаружился еще кот, черный в белых носочках, Кузя. Шерсть у него немного светится на шее и спине, но не лишай, оказывается. Дерматит. Кота возили к ветеринару, делали уколы.
    
    О бытовых условиях. Отопление центральное и печное, в окнах стеклопакеты. Жарко, на самом деле, почти как в поезде. Туалеты чистые, вода из-под крана достаточно вкусная, смена постельного белья раз в две недели. Баня по четвергам и пятницам.
    
    В местах общего пользования (коридоры, лестницы, гардеробы) свет включается датчиками движения. В туалете, когда включен свет в кабинке, работает вытяжка.
    
    \mds{2017-01-30, понедельник}. До обеда~---~развозка дров по корпусам, для чего имеются сани, запряженные мерином по кличке Иртыш. Работа не слишком тяжелая, Иртыш славный. После обеда~---~колоть дрова, это уже требует сноровки, которой у меня нет. Под конец немного стало получаться.
    
    Завтра обещают кухонный наряд, он же послушание. Дрова~---~тоже послушание.
    
    \mds{2017-01-31, вторник}. На кухню заступать надо в шесть утра, зато есть большие перерывы в течение дня. Очень удобно, поскольку от кельи до кухни не более двух минут пешком. Тут вообще самый дальний поход не более пяти минут. Удивительно экономит жизненные силы. Занимаюсь накрытием столов в трапезной. Картошку чистить, видимо, не дорос еще.
    
    В процессе накрытия есть один драматический момент: вдвоем донести фердинанд с горячим борщом от плиты до грузового лифта на кухне, а потом до служебного стола в трапезной. Это не запредельно тяжело, но цена оплошности очень велика: и обваришься весьма основательно, и оставишь монастырь без обеда.
    
    Завтра опять кухонное послушание.
    
    \section*{февраль}
    
    \mds{2017-02-01, среда}. К завтраку забыл на некоторых столах дополнить комплект тарелок, а вставать из-за стола нельзя никому... Неловко получилось, хотя никто сильно не обиделся.
    
    \mds{2017-02-02, четверг}. Сегодня была жареная рыба: щука, судя по головам, оставшимся на кухне. До того, чаще из консервных банок. Еще случился банный день, как нельзя кстати. Баня небольшая и чистенькая, на четырех человек. Выглядит хорошо как изнутри, так и снаружи. Можно попариться и помыться из шайки. Ни спешить, ни экономить воду не надо. Одежду стирают централизовано, нижнее белье каждый стирает сам.
    
    \mds{2017-02-03, пятница}. К обеду ждали игумена о. Иосифа, не дождались.
    
    Второй день уже минус двадцать. В келье по-прежнему жарко, на улице вполне терпимо, даже хорошо: ветра нет и сухой воздух. Коты, включая Кузю с дерматитом и редкой местами шерстью, чувствуют себя на снегу хорошо, и помещение покидают по собственной воле. Из братского корпуса, где моя келья, в Казанский собор, где трапезная, хожу в не слишком теплых носках и резиновых тапочках: холода ощутить не успеваешь совершенно, зато ноги не преют во время работы в помещении.
    
    Завтра буду снимать с веревки выстиранную одежду. Посмотрим, как она сохнет в таких условиях.
    
    \mds{2017-02-04, суббота}. Одежда на морозе таки сохнет не очень хорошо, колом стоит. Отойдя в помещении, оказывается, все же, не мокрой, а влажной, досушить нетрудно.
    
    Вечером, после работы в трапезной, пришлось еще мыть полы на кухне. К счастью, не слишком грязные.
    
    Завтра только две трапезы, обед в десять утра и ужин в обычное время. Выход на работу зачем-то в семь утра. Обед перенесен на девять, оказывается.
    
    Вообще, монастырь действительно похож на подводную лодку: в любое время суток кто-то спит, а кто-то бодрствует.
    
    \mds{2017-02-05, воскресенье}. Минус тридцать. Резиновые тапки перестают скользить на утоптанном снегу при минус двадцати. Иными словами, при минус десяти поскользнуться в тапках не просто, а очень просто. Вот и сегодня ходить между братским корпусом и Казанским собором очень удобно.
    
    \mds{2017-02-06, понедельник}. Минус сорок два. Звездное небо, в шесть утра, на зависть Канту. В келье по-прежнему жарковато, хотя, может, и полегче, чем при минус пяти. Отопление, с печами и батареями, фундаментальное, а стены если и меньше метра в толщину, то не намного.
    
    Завтра обещают новое послушание. Какое, неясно. В знак же окончания этого состоялось мытье полов в трапезной и далеко не только, работа на двух человек. Управились к часу ночи, итого восемнадцать часов разнообразной деятельности подряд.
    
    \mds{2017-02-07, вторник}. Проснулся в четверть седьмого. За завтраком выяснилось, что послушание знакомое: дрова. Удивительно, сколько рук и мест всяких проходят они на пути из леса к печке. Просто золотыми должны быть дрова эти. Думается, стоили они дороже, чем нам отсюда кажется, во времена якобы идиллические, когда, кроме дров, ничего не было.
    
    Пошел снежок, и температура поднялась до минус пятнадцати.
    
    Собаки тут, все же, какие-то есть, но далеко на периферии, где бревна на дрова лежат, и совсем не чувствуют себя хозяевами.
    
    Познакомился с о. Феофаном. Нашелся, наконец, человек, понимающий компьютерные слова.
    
    \mds{2017-02-08, среда}. Снова дрова, и уже совсем тепло, около нуля.
    
    Завтра баня, в половине пятого.
    
    \mds{2017-02-09, четверг}. По-прежнему тепло, а с дров перевели на рыбу чистить. Посмотрим, каково оно.
    
    Ничего страшного, эта рыба путассу. Управились до обеда, и теперь только баня, т.е. сплошные удовольствия.
    
    Любопытно: температура не выше нуля, и второй день только, а на карнизах довольно серьезные сосульки. Отопление чувствуется?
    
    После бани поменял белье. Особых формальностей нет: дали простыню, пододеяльник и наволочку, полотенца менять не стал, а лишние трусы получил без вопросов.
    
    \mds{2017-02-10, пятница}. Подушка, в отличие от постельного белья, оказалась проблемой. Заменить ее можно только с благословения о. Рафаила. Подушка велика, наволочки толком не налезают, но это полбеды: из нее перья лезут, а это не радует не только меня.
    
    С помощью о. Матфея записался на прием в местную поликлинику, по поводу давления. Где гипертонию лечить, как не здесь. Заодно посмотрю на окрестности.
    
    \mds{2017-02-11, суббота}. Сегодня была рыба, четырех видов: горбуша, минтай, щука и окунь. Легче всего чистится горбуша, хуже всего~---~окунь. Этого последнего хорошо бы занести в Красную книгу, ко всеобщей пользе и удовольствию. Есть нечего, а чистится как акула.
    
    Наконец-то завезли гречку и сварили из нее кашу, ко всеобщему удовольствию.
    
    \mds{2017-02-12, воскресенье}. Завтра еду в поликлинику. Не один~---~в составе экспедиции. Вечером испросил благословения о. Иосифа, как положено. Без благословения игумена территорию монастыря не покидают. Показалось, что о. Иосиф остался слегка удивлен.
    
    \mds{2017-02-13, понедельник}. В девять утра выехали в Карпогоры на джипе, в одиннадцать четырнадцать прием у терапевта (Колик Сергей Николаевич). Судя по виду, про гипертонию он должен знать все, из первых рук.
    
    Померив давление, сходу отправил меня на десятидневный курс в дневном стационаре. Дневным стационар называется потому, что в нем не ночуют, лечатся же с восьми утра до двух часов дня. Когда я все это понял, опытные мои попутчики, включая о. Матфея, куда-то потерялись вместе с автомобилем. Впрочем, когда я уже собрался искать подворье самостоятельно, о. Матфей, правильно рассудивший, что дальше поликлиники и стационара пациент никуда не денется, меня нашел.
    
    Чрезвычайно любезная Галина Никифоровна, оформив мне временную медицинскую страховку, отправила заявление на постоянную, которая должна поступить в течение месяца или около того. Ей вполне хватило моих устных объяснений насчет бывшей бессрочной, но устаревшей страховки, и номера СНИЛС, продиктованного по телефону.
    
    Словом, ничто не помешало поставить мне капельницу с аспартамом и сделать инъекцию никотиновой кислоты внутримышечно еще до обеда.
    
    Купил, по необходимости, мегафоновскую симку, только для разговоров, без интернета. Тарифы, выложенные для покупателя, даже не смешны: триста рублей в месяц за полтора гигабайта интернета в месяц же. В качестве предмета роскоши (VIP тариф) предлагается пятнадцать гигабайт в месяц за две тысячи рублей. Мегафон, действительно, странная организация. Покупателю также предлагаются зарядные устройства за восемьсот рублей, и наушники-пробки за четыреста. Еще бабушке при мне возвращали две тысячи рублей за подписки, как бы сделанные ею при помощи кнопочного мобильника минимальной конфигурации. Теперь я понимаю, почему о. Феофан предложил мне первым делом, по возвращении, явиться с мегафоновской симкой к нему, для вычесывания из этого жульнического изделия блох. Посмотрим, что ему удастся сделать.
    
    На подворье меня поселили в келью на шесть человек, одного. Храпеть можно в свое удовольствие. В келье не холодно и не жарко, в отличие от Верколы. Есть письменный стол, полка с книгами.
    
    Кормят и здесь хорошо, туалет просторный, с душевой кабинкой.
    
    Есть шикарная трехцветная кошка четырнадцати лет, зовут Муся.
    
    \mds{2017-02-14, вторник}. Кровь на анализ, капельница, укол. На подворье все вкусно и замечательно. Чай на выбор, хоть травяной, хоть из пакетика.
    
    \mds{2017-02-15, среда}. Галина Никифоровна утверждает, что новый полис для меня был заказан в декабре 2016 года. Кем, интересно? По счастью, проблемы это не составляет. У Галины Никифоровны проблем, почему-то, не возникает вовсе.

    \mds{2017-02-16, четверг}. По дороге в стационар и обратно заглядывал в магазины. Платить картой можно практически везде. На почте, помимо стандартных услуг, торгуют подсолнечным маслом, макаронами, рыбой и тушенкой.
    
    Мелочевку, вроде флешек и наушников, надо смотреть во <<Флагмане>>, а не в <<Березке>>, где офис Мегафона. Подумалось, что отцы-основатели <<Березки>> уже не в курсе значения этого слова применительно к магазинам.
    
    О. Анатолий выдал мне <<Букварь школьника>>, где в начале есть краткий учебник церковнославянского языка А.В. Ушкова, около сотни страниц. <<Букварь>> целиком напоминает, по объему, полное собрание сочинений Шекспира с комментариями.

    \mds{2017-02-17, пятница}. Сдал всякие excreta на анализ. Набор контейнеров обошелся в пятьдесят рублей, что показалось недорого, на фоне прочих карпогорских цен. Дорого тут все, никак не дешевле, чем в большом городе.
    
    Еще нашел в бумажках направление на УЗИ. Хорошо, что сейчас, а не потом.

    \mds{2017-02-18, суббота}. Был на дневной службе в девять утра. Прихожан человек тридцать, из них трое мужчин. Женщин, не достигших менопаузы, тоже почти нет. Гена Магилюк звонил из Краснодара, хорошо ему: казенная мегафоновская симка~---~триста рублей в месяц, объем разговоров по стране не ограничен. Могут, когда хочут.
    
    Сегодня меня долго убеждали в том, что вчера я уже достал из погреба картошку. Моей глухой несознанке верить не хотят.

    \mds{2017-02-19, воскресенье}. Дошло дело и до неиллюзорной картошки, наконец. \emph{Втроем}, не без усилий, достали три мешка картошки общим весом не менее ста пятидесяти килограммов.
    
    Погреб вырыт в грунте, слегка на отшибе, прикрыт сарайчиком. Вход представляет собой круглое отверстие в земле, облагороженное большой покрышкой; туда еще надо спустить лестницу. Над отверстием одиночный блок с веревкой, толстой и мягкой, к счастью. Мешок, весом никак не меньше пятидесяти килограммов, берется на крюк; поднимать его нелегко еще и потому, что, кроме физической силы, требуется приличный вес для сцепления с грунтом. Три человека, один в погребе, два на блоке, самое то. Вдвоем, теоретически, можно, но мужики нужны \emph{очень} крепкие. А я позавчера был реальный супермен.
    
    Непонятной выдумке, оказывается, может хватить двадцати четырех часов, чтобы окрепнуть в почти правду.

    \mds{2017-02-20, понедельник}. УЗИ двадцатого и двадцать первого не работает: выезд в Верколу. Во-вторых, записи на февраль уже нет. В-третьих, отпуск у них, с первого марта. Надо, значит, не пропустить следующего выезда в Верколу, не раньше апреля, как я понимаю.

    \mds{2017-02-21, вторник}. Пропустил сегодня стационар из-за мороза. Вышел, и сразу повернул обратно. Не стоит аспартам воспаления легких.
    
    Днем гораздо теплее, но идти поздно уже.

    \mds{2017-02-22, среда}. Звонил Гена Магилюк. Рассказал, что в Пинеге водятся, в основном, щука и окунь. Очень похоже на правду, судя по нашему столу. Окунь, что хорошо видно по икре и молокам, достигает половой зрелости задолго до достижения серьезных размеров. Водится, опять же, повсеместно. Не будет мне Красной книги для окуней.
    
    Поймал еще в стационаре за пуговицу терапевта своего, Колика Сергея Николаевича. Говорит, выпишет меня стационар сам, без его непосредственного участия.
    
    Далее четыре выходных, в том числе и для медицины. Дырки от уколов подживут, и это неплохо, но пребывание мое в Карпогорах опять затягивается.

    \mds{2017-02-23, четверг}. Вернулся о. Анатолий из Архангельска, привез, среди прочего, несколько пластиковых бутылей с лампадным маслом. Масло белое, мутное~---~не иначе, замерзла на холоде водяная суспензия. Так и не спросил, что лампадное масло собой представляет.
    
    Живут на подворье еще две собаки, в вольерах из досок. Немецкие овчарки сомнительного происхождения, судя по экстерьеру. Лают и кидаются, в том числе с разбега, если на свободе, но на куски не рвут.

    \mds{2017-02-24, пятница}. Собак зовут Окса и Рика, сестры шести лет. Отличить легко: у Оксы не стоят уши. Крепкие, быстрые и тяжелые, но гораздо дружелюбнее, чем кажутся.
    
    Выпустил их из вольеров на свой страх и риск, собаки очень обрадовались. Сидят практически сутки напролет без движения, жалко.

    \mds{2017-02-25, суббота}. Для разнообразия, подъехал трактор сгребать снег. Облегчение, конечно, но лучше бы подъехал завтра: снега, судя по всему, будет больше, чем сегодня.

    \mds{2017-02-26, воскресенье}. Так и вышло. Снег шел всю ночь, но убирать его будем вручную.
    
    Звонили из Пскова. Слепую кошку оперировали по поводу опухоли где-то в ухе. Теперь она и слышит одним ухом только. Остатки глазных яблок тоже надо удалять.

    \mds{2017-02-27, понедельник}. Выбивали ковры. Снег под ними черный. Всегда удивлялся, как хорошо они пыль собирают.
    
    На прощание о. Анатолий подарил мне <<Краткий учебник церковнославянского языка>> А.В. Ушкова, очень кстати. Это начало <<Букваря школьника>>, изданное отдельно. Около ста пятидесяти страниц, формат небольшой, удобно.

    \mds{2017-02-28, вторник}. Вот и последний день в Карпогорах. Выписали <<с улучшением>>, выдали ксерокопию выписного эпикриза. Почерк лучше стандартного врачебного, но все равно понимаю плохо. Ничего, о. Матфей разберется.
    
    Пообедал в Карпогорах, ужинал в Верколе. Самой езды полчаса, если ехать быстро и нигде не останавливаться.
    
    \section*{март}
    
    \mds{2017-03-01, среда}. В шесть утра впервые заметил, что на печах в коридоре братского корпуса присутствует трехлучевой орнамент, который triskelion. Всегда считал его кельтской древностью, но словарь Random House относит термин к середине XIX века, а сам логотип к Сицилии и острову Мэн. Что, с другой стороны, кельтскому происхождению не противоречит.
    
    Печи новые, явно не времен основания монастыря.
    
    Послушание сегодня действительно пыльное, в первый раз: печь разбирать. Со склада подходящую одежду просто так не возьмешь, надо ловить кладовщика. Взять заранее тоже не выходит~---~негде хранить. Я пробовал, собственно. Приходится работать в относительно чистом, полагаясь на послезавтрашнюю баню и возможность сдать в стирку сразу все. Нехорошо как-то. На складе именно того, что надо, просто некуда девать, а народ работает в своем.
    
    О. Матфей выдал из своих закромов выписанные мне лекарства. Надо почитать вкладыши.
    
    Еще один кот нашелся, старожил. Десять лет ему, зовут Чижик. Живет на втором этаже братского корпуса, безвылазно, по-видимому, иначе я бы его раньше заметил. Выглядит печальным и усталым. Черно-белый, вроде Кузи, но у того кончик хвоста белый.
    
    \mds{2017-03-02, четверг}. Сегодня снова дрова, дело уже привычное.
    
    Ездил после обеда с о. Матфеем на почту, где он получил посылки на имя о. Рафаила, эконома, которых как раз хватило, чтобы заполнить багажник <<Шкоды Октавии>>. Получил также благословение о. Рафаила на вторую наволочку. В первой лежит сложенное одеяло, которое удобно ночью подкладывать под голову: меньше храплю. Вторая нужна для собственно подушки, вместе с которой удобно делать высокое изголовье, если надо читать-писать.
    
    По прямому назначению одеяло употреблять не приходится~---~слишком жарко.
    
    \mds{2017-03-03, пятница}. Принял сегодня за завтраком первую таблетку индапамида. Если верить вкладышу, заметного эффекта не следует ожидать раньше, чем через четыре недели. Посмотрим.
    
    О дровах, тонкости. Нетрудно догадаться самостоятельно, что сухие дрова предпочтительнее сырых. Но есть нюансы: сырые дрова хвойных пород в печку еще можно, а сырые березовые уже нельзя. Такие дела.
     
    Баня в четверть четвертого. Оставил в стирку все-все-все, включая коричневую куртку. Последствия разборки печи в игуменском корпусе.
    
    \mds{2017-03-04, суббота}. Для работы на дровах пришлось, за недоступностью своей, получать казенную куртку. Выдали камуфляжный бушлат, тяжелый и брутальный. Не сильно поношенный, впрочем, и добротный.
    
    О. Феофан завел личный кабинет для моей мегафоновской симки, подобрал вменяемый тариф, убил подписки и вообще навел порядок. Спасибо ему, благодетелю. Осталось положить денег на телефон, что не так просто, как можно подумать: сбербанковская карта завязана на Tele2, а он ловится не ближе Карпогор. От наличных какой-то толк бывает, по слухам, на почте, но до нее тоже добраться надо.
    
    Светлая сторона: андроидный телефон Арч видит в качестве сетевой карты прямо из коробки. В отличие от Мака.
    
    \mds{2017-03-05, воскресенье}. О. Феофан завершил воспитание мегафоновского интернета на моем телефоне. Обновлений системы ноутбука восемьсот метров, процесс пошел около пяти вечера. Интернет, само собой, тот еще, обновления накапливать нельзя.
    
    \mds{2017-03-06, понедельник}. Шесть утра. Обновления еще не вполне скачаны, хотя осталось всего семнадцать метров.
    
    Семь утра. Первый коммит на новом месте, ага.
    
    Половина восьмого. Скачано все из репозиториев.
    
    Восемь утра. Обновление завершено, кроме AUR. А теперь дрова, дрова, дрова.
    
    \mds{2017-03-07, вторник}. Разгрузка камаза с предметами снабжения и дрова. В числе предметов снабжения оказались пятидесятикилограммовые мешки с цементом и шамотный кирпич. Но, в целом, рабочий день закончился удачно, за два часа до обеда. Даже немного раньше.
    
    \mds{2017-03-08, среда}. Дрова. Наделал фоток дровяного хозяйства, пользуясь ярким солнцем.
    
    \mds{2017-03-09, четверг}. Дрова, много. Завтра баня, так что любая работа только до обеда.
    
    \mds{2017-03-10, пятница}. Поменял постельное белье, сдал пару пустяков в стирку.
    
    \mds{2017-03-11, суббота}. Дрова вчера, сегодня, но уже не завтра. Грядущее воскресенье очень кстати.
    
    К вопросу о переправе через Пинегу: сегодня УАЗ-452 таки попортил лед, хоть и не утонул. Притом дорожные знаки, разрешающие переправу, еще не сняты. Камазам, стало быть, нельзя уже совсем; людям пока можно.
    
    \mds{2017-03-12, воскресенье}. Получил на складе бутылку минеральной воды <<Куртяевская>>. Здоровье из недр поморской земли, как и сказано на этикетке.
    
    Начиная с понедельника, о. Матфей будет мерить мне давление. Каждый день, возможно.
    
    \mds{2017-03-13, понедельник}. Шесть утра. Рыжий сидит на том же месте, с которого был согнан вчера, буквально пинками (лестница на второй этаж братского корпуса). Оказывается, лестница эта принадлежит Моне, который черный и пушистый. Они как раз начали было разбираться по-кошачьи, кто тут главный, но вмешался человек, на стороне Мони. Стоило трудов?
    
    \mds{2017-03-14, вторник}. Ходили пешком за реку, грузить кирпичи на тракторный прицеп, потом разгружали рядом с дровами, и так два раза. На льду,---~точнее, на утрамбованном снегу переправы,~---~появились трещины. Узкие, палец не просунешь, но извилистые и длинные. Впрочем, трактор с кирпичами под лед и не думал проваливаться, так что субботняя история про уазик мне не ясна.
    
    \mds{2017-03-15, среда}. Опять кирпичи. Теперь погрузка обратно в прицеп, и снова разгрузка, уже игуменский корпус. По счастью, первый этаж. И в таких тепличных условиях, нагрузка на человека~---~три кирпича. Три, больше~---~чревато. А кто-то все это строил, стены в метр толщиной, собор. Из кирпичей.
    
    Дров, как работы, хватит от силы недели на две, потом в сарае просто не будет места. Что дальше, интересно?
    
    \mds{2017-03-16, четверг}. Жареные кальмары очень хороши, и никакого послевкусия.
    
    Будем сегодняшний день считать началом работы над версткой учебника Ушкова. Двадцать четвертая страница пошла, но до того, на самом деле, была подготовка такая и подготовка сякая, незаменимый интернет туда же.
    
    \mds{2017-03-17, пятница}. На кухне говорят, жареные кальмары оказались большой ошибкой. Исчезли, никто глазом моргнуть не успел.
    
    Баня хороша, как всегда. Жаль только, что раз в неделю. Впрочем, куда больше~---~по факту, наполовину выходной день.
    
    Говорят, грачи прилетели. Кто-то где-то видел. Зима, стало быть, уже не вернется. Не должна.
    
    \mds{2017-03-18, суббота}. Длинный день получился. Сделал шесть страниц учебника, не считая дров, куда же без них.
    
    Вечером пообщался с Эдиком по WhatsApp.
    
    \mds{2017-03-19, воскресенье}. На кухне исправили давешнюю ошибку: вместо жареных кальмаров теперь котлеты из кальмаров же. Далеко не то, но уходят все равно неплохо.
    
    Четырех человек отправили на расчистку пасеки. Для воскресенья большая редкость.
    
    Говорят, будто теперь Tele2 в монастыре ловится. Или будет ловиться? Мой индикатор ничего не показывает.
    
    \mds{2017-03-20, понедельник}. Простыл по весне. Вчера еще, но и сегодня рано судить, насколько.
    
    Со склада выдали орехов, лесных и грецких, и еще фиников, на келью. Вчера выдавали сушки и конфеты, не считая минеральной воды. Процветаем, однако.
    
    История с Tele2 напоминает историю с карпогорской картошкой. Ну нету Tele2 никакого, ни один индикатор ни на одном телефоне ничего не показывает, но идея живет.
    
    \mds{2017-03-21, вторник}. Для памяти: сегодня убирался в келье~---~пылесос, тряпка, все дела. Уборка у нас чаще и тщательнее среднего, через три дня на четвертый, по очереди. Чем больше в келье народу, тем реже приходит очередь; сейчас нас трое.
    
    Грибной суп дивный: грибов, хоть и не белых, там больше, чем картошки.
    
    \mds{2017-03-22, среда}. Ремонт оконных проемов в келье. Пылища, неподъемные мешки со строительным мусором. Завтра, предположительно, только мокрая работа. Все, что вам надо знать про очередность уборки.
    
    \mds{2017-03-23, четверг}. Нет, ремонт продлится дольше.
    
    \mds{2017-03-24, пятница}. Баня. Удалось договориться насчет покрывал~---~их, вообще-то, не принимают в стирку, как и нижнее белье. Иначе пыль в них так и осталась бы.
    
    Утром чистил кальмаров. Просто прелесть, с окунями никакого сравнения.
    
    \mds{2017-03-25, суббота}. Работы на дровах уже почти нет, одна развозка. Хорошего в этом мало: разбирать строительный мусор на чердаке игуменского корпуса куда как неприятней.
    
    Ремонт оконных проемов подходит к концу. Осталась уборка~---~ну, и побелка.
    
    \mds{2017-03-26, воскресенье}. Все спят. Служба с восьми утра до половины десятого, не длинная. А потом все спят, просыпаясь только на обед и на ужин. Завтрака по воскресеньям нет.
    
    \mds{2017-03-27, понедельник}. Разборка чердака и части перекрытий в игуменском корпусе. В окно, для последующего вывоза, выбрасываются кирпичи, штукатурка и отслужившие свое деревянные балки, в огромном количестве. Пыль удушающая, на самых трудных участках работают в респираторах. Поразил кот Моня, явившийся в пыльное царство пообщаться. Когда кот ошивается на кухне, это понятно. Моня же систематически посещает любые скопления людей.
    
    \mds{2017-03-28, вторник}. Вернули меня на кухню. Как всегда, неожиданно. Ничего, пока без приключений. Внезапно назначили УЗИ на завтра, в 14:15. Вот как бывает. Не надо ждать неизвестно сколько.
    
    \mds{2017-03-29, среда}. В Карпогоры ехал вместе с настоятелем, о. Иосифом. Вышел конфуз: надо было сразу идти на другой берег Пинеги, а я не понял. Не зря военных учат составлению приказов. Сказал бы мне кто "в 07:45 выдвигаешься на дальний берег р. Пинеги, где тебя будет ждать машина такого-то цвета с таким-то номером"~---~и что тут можно перепутать?
    
    УЗИ продолжалось не слишком долго~---~ровно столько, сколько надо, чтобы опоздать на машину настоятеля, идущую обратно в Верколу. Теперь ночевать на подворье.

    \mds{2017-03-30, четверг}. На подворье сплю в той келье и на той же постели, которую оставил месяц назад, двадцать восьмого февраля. Ровная до чего жизнь.

    \mds{2017-03-31, пятница}. Вернулся из Карпогор с о. Иосифом, ближе к вечеру. Баню, соответственно, пропустил, удалось только помыть голову под карпогорским душем, что совсем не просто: отрегулировать температуру воды надо уметь. Наладил мне душ Николай, которому я убил хорошо спрятанные платные сервисы на мегафоновском номере (УСТЗАПРЕТ1 и НЕТКЛИК1, SMS по адресу 5051).
    
    \section*{апрель}

    \mds{2017-04-01, суббота}. Сервировка столов, полет нормальный. Закончился, наконец, ремонт в келье. Совсем закончился, с мытьем и уборкой. Подходил к концу, если проследить, неделю назад, в субботу, двадцать пятого марта. Начался, соответственно, в среду, двадцать второго. Итого, не менее десяти дней. Вместо двух по плану.
    
    \mds{2017-04-02, воскресенье}. Трапезная, котлеты из кальмара, салат из свежей капусты.
    
    Вернул о. Иосифу деньги, данные мне на дорогу из Карпогор, которые не пришлось истратить.
    
    \mds{2017-04-03, понедельник}. Мытье полов по окончании трапезного послушания. Закончили где-то в половине второго ночи. Отмыли декоративную плитку в трапезной, всего-то две узких полосочки на полу. Дизайнеры явно не думали, что делают.
    
    \mds{2017-04-04, вторник}. Отказался от завтраков. Манная каша пополам с малиновым вареньем не стоит этого беспокойства. Не Буратино я. Монахи вон тоже не завтракают.
    
    Простуда, впервые замеченная в воскресенье, девятнадцатого марта, окончательно не проходит, хоть и не усугубляется. Надо бы полежать все-таки.
    
    \mds{2017-04-05, среда}. Первый день проведен в качестве болящего, официально. Посмотрим, как оно лежится в келье. От Гены таки получил баночку малинового варенья в качестве лекарства.
    
    \mds{2017-04-06, четверг}. Сходил в баню. Постельное белье поменять не смог, так как сегодня, по случаю предпраздничного дня, выдавали его не в семь, как обычно, а в пять, чего я не знал. На обеде объявляли~---~не больно-то хороша там акустика, до соседнего стола мало что доходит, даже когда вслушиваешься специально.
    
    \mds{2017-04-07, пятница}. Сопли текут ручьем, но настроение бодрое.
    
    \mds{2017-04-08, суббота}. Удивительно, но от учебника Ушкова осталось не более десятка страниц работы по вводу. Потом и учить можно будет никуда не торопясь уже.
    
    Бытует мнение, что кальмарами нас кормили по случаю поста, так как не рыба и не мясо. Что гады морские, это не страшно. Мнение народное, настоятеля спрашивать не берусь.
    
    \mds{2017-04-09, воскресенье}. Закончил ввод текста учебника церковнославянского языка (Ушков). Условное начало работы~---~шестнадцатое марта. До того подготовительные мероприятия, трудно судить; но до шестого марта (даты подключения интернета) никакого старта и быть не могло, такие дела. Можно смело сказать, что уложился в четыре недели. Сто пятнадцать страниц достаточно нетривиального текста.
    
    За этот месяц понемногу выяснилось, что в монастыре учебник никому особо не нужен. Ничего, в сетевую эпоху не страшно.
    
    \mds{2017-04-10, понедельник}. Официально объявлен выздоровевшим; завтра опять в трапезную. Разобрался, как называется универсальный головной убор (buff). Остальное просто.
    
    \mds{2017-04-11, вторник}. На кухне, при мойке и в других местах не осталось ни одного неопытного человека. Очень удобно. Если бы еще не генеральная уборка, которой все это счастье кончится\ldots{} Кто угодно может сказать, и скажет обязательно, что <<вы только грязь размазали>>. Ну да, швабра размазывает грязь, это ее предназначение. Если палубу корабля захлестывают волны, размазанная грязь куда-то уходит, а так не очень. Никуда не денешься.
    
    Колокол, в который бьют перед трапезой, висит над широким поручнем у входа в братский корпус. Сегодня Моня лежал на поручне аккурат под самым колоколом, это около полуметра от источника звука, \emph{очень} громкого. В момент удара он и не подумал никуда бежать; если чем и был слегка недоволен, так только мельтешением над головой. Так лежать и остался. Удивительно стойкое животное.
    
    \mds{2017-04-12, среда}. Баня, неожиданно, сегодня. Мог бы вообще пропустить. Постельное белье поменять снова не вышло: освободился из трапезной в половине восьмого вечера. На реальную неприятность, однако, еще не тянет.
    
    \mds{2017-04-13, четверг}. Коровам (или курам?) достается грибной суп, в изрядном количестве.
    
    \mds{2017-04-14, пятница}. Длинный день, а завтра обещают длиннее, хотя завтрака не будет.
    
    \mds{2017-04-15, суббота}. Обед и ужин прошли, как обычно. Всю ночь (воскресенье) еще разговляться. Интересно, что будет в воскресенье днем?
    
    \mds{2017-04-16, воскресенье}. После ночного напряжения нет, хотя бы, завтрака. Обед и ужин прошли все-таки легче ночной трапезы, по накатанной колее. Праздник равен пожару.
    
    \mds{2017-04-17, понедельник}. Снова не было завтрака, зато сдача смены и генеральная уборка. Освободился в три часа ночи, чуть живой.
    
    \mds{2017-04-18, вторник}. Выходной, наконец. Выбрался из кельи только на обед. На ужин и ходить не стал.
    
    В монастырской жизни намечаются сезонные изменения: разъезжаются трудники. Зимой многих держит здесь безработица, летом же больше работы, если на стройке, например. Кроме того, приближается ледоход на Пинеге, когда выбраться из монастыря станет практически невозможно. Уже уехало восемь человек, самых трудоспособных.
    
    На время ледохода останавливается и почта.
    
    \mds{2017-04-19, среда}. Немножко дров до обеда, и жизнь продолжается. Моня очень хотел помочь, даже запрыгивал на сани, куда дрова в это самое время \emph{укладывались}. Сделал фотку, как Моня сидит на чурбане на фоне дров, но степени активного вмешательства в дела человеческие на ней совсем не видно.
    
    \mds{2017-04-20, четверг}. До обеда ни завтрака, ни послушаний.
    
    \mds{2017-04-21, пятница}. Получив благословение о. Иосифа, сходил на почту за пришедшей вчера посылкой от Галины, чтобы не лежать ей там до окончания ледохода. Благословение дано было в форме необычной, <<на свой страх и риск>>. Рассудив, что при действительно ненадежном льду благословения не дадут никакого, таки пошел. Лед, конечно, держит пока человека вполне надежно, даже на лужах в дорожной колее не трещит. 
    
    Полюбовался на аппарат на воздушной подушке, припаркованный на берегу Пинеги. Вид вполне заводской, но никаких логотипов производителя найти не удалось. Он не просто так стоит: во время ледохода никакого другого способа пересечь Пинегу не существует. 
    
    \mds{2017-04-22, суббота}. Немножко дров, на которые я даже не успел, немножко песка в мешках. Работают три человека: один копает, другой мешок держит, я ношу. Правда, только ношу и сбрасываю, поднимают его мне на спину другие.
    
    \mds{2017-04-23, воскресенье}. Настоящий выходной. На обед котлеты из горбуши и сливочное масло в изобилии.

    \mds{2017-04-24, понедельник}. Укладываем красиво доски, оставшиеся после разборки покосившегося сарая.
    
    \mds{2017-04-25, вторник}. Чищу и режу горбушу на кухне. Чистка ничего, а резать сложно. Ножи тупые, шкуру берут особенно плохо, приходится ножницами.
    
    \mds{2017-04-26, среда}. Возим песок в мешках с берегов Пинеги. Песок идет в раствор, раствор в пол первого этажа братского корпуса.
    
    \mds{2017-04-27, четверг}. Продолжается уборка в игуменском корпусе. Больше не такая пыльная, как в начале. Интернет уж недели две никакой.
    
    В игуменском корпусе пожар, оказывается, был. Не такой, чтоб выгорела крыша со стропилами, сплошное дерево, а такой, что выгорели перекрытия ниже и \emph{печи}. Да, именно поэтому их сейчас и перекладывали. Детали я понял плохо, хотя и переспрашивал.
    
    \mds{2017-04-28, пятница}. Устройство кельи для паломников в игуменском корпусе. Койки старые, сетка и в лучшие времена была не очень, оттого покрываются деревянными щитами. Пылесос, ведро, тряпка. Немного и не тяжело.
    
    В четверть четвертого баня.
    
    \mds{2017-04-29, суббота}. Рыба, которой меня теперь не испугаешь, если это не окунь. Горбуша, щука и таки окунь, но его теперь не надо чистить. Никого не надо чистить, надо делать филе. Хребет, по возможности с ребрами, извлекается. Рыба при этом распластывается, а мясо соскребается со шкуры. Куда как менее кровопролитный процесс, чем многократное перерезание хребта тупым ножом. Острых нету, да. Окунь же пошел на сковородку прямо в чешуе. Кто-то накапливает опыт, не иначе.
    
    Еще чистил щучьи головы. Удалял жабры, то есть. Головы потом варят, и подают отдельно, как деликатес. Никто особо не ест, впрочем.
    
    \mds{2017-04-30, воскресенье}. Все спокойно, как и полагается выходному. На обед, предсказуемо, котлеты из вчерашнего филе.
    
    \mds{2017-05-01, понедельник}. Развозим дрова. Вместо саней уже телега с колесами от жигулей. После обеда очистка коек в игуменском корпусе от пыли.
    
    \mds{2017-05-02, вторник}. Перенесли посевную картошку из погреба в коридор первого этажа игуменского корпуса. Делается это каждый год в начале ледохода. Итого, между последним моим пешим переходом через Пинегу и реальным движением льда~---~менее двух недель.
    
    Моем столы и пол в келье на втором этаже гостиничного корпуса, где была швейная мастерская.
    
    \mds{2017-05-03, среда}. Заделываем потолок в той же келье на втором этаже гостиничного корпуса. С последующим мытьем столов и пола.
    
    \mds{2017-05-04, четверг}. Мыть пол в третий раз так и не пришлось: перебросили на элитную посевную картошку, которую надо выгружать из погреба, раскладывая по деревянным ящикам. Правда, картошка эта ничем с виду не отличается от обыкновенной. Есть опасения, что перепутают ее в конце концов, благо раскладываем по таким же ящикам, и стоять те ящики будут на том же полу в том же игуменском корпусе.
    
    \mds{2017-05-05, пятница}. Продолжается выгрузка элитной посевной картошки из погреба. Пол в швейной так и остался подметенным только.
    
    \mds{2017-05-06, суббота}. Попал на чистку картошки, наконец-то. Вчетвером достаточно быстро и весело. Картошка хорошая и плотная, только много мелкой.
    
    \mds{2017-05-07, воскресенье}. Много жареной горбуши и упрямого Vim'а. В остальном спокойное воскресенье, как и полагается.
     
    \mds{2017-05-08, понедельник}. Уложено в ящики около 1300кг посевной картошки, элитной и попроще. Этого недостаточно. В погребе открыли летучую мышь. Прямо под сводом, у стены, висела на виду все это время. Мышь спит, никто не замечает.
    
    \mds{2017-05-09, вторник}. Картошку укладываем, ящики носим. Нужно около двух тонн, при расчетном урожае двенадцать тонн. Ящиков, без сомнения, не хватит. Десятикилограммовый ящик выматывает на удивление хорошо: никак не думал, что там \emph{всего} десять килограммов. Носить, правда, приходится где внаклонку, а где по лестнице, но все же.
    
    \mds{2017-05-10, среда}. Заполнили остаток ящиков, оставшуюся посевную картошку разложили ровным слоем в отсеках погреба. Она там, будто бы, во благовремении прорастет. Посмотрим. Завтра опять туда, хотя характер работы пока не ясен.
    
    Моня сегодня заглянул к нам в келью, где умылся и поспал немного на свободной койке.
    
    \mds{2017-05-11, четверг}. Освобождение от старых кроватей тупичка в игуменском корпусе. Лестница, площадка, ничего особенного. Очистил от содержимого пылесос <<Кархер>>~---~солидная машина. Потом баня, вместо пятницы. За ужином поставили читать. В первый раз, и едва ли в последний.

    \mds{2017-05-12, пятница}. Дрова, в том числе колоть. Старые, сорные: остатки от снесенных построек. Чем трухлявее, тем легче работается.

    Главное событие для: в келье паника. Подселяют вновь прибывшего, на ночь глядя. С поезда видимо, как и меня в свое время. Только меня подселили вовсе без предупреждения.

    В 23:30 выходить на крыльцо в сапогах. Идти на берег, за вещами. Жизнь делается интереснее с каждым часом.

    В половине двенадцатого <<лодочнику еще не звонили>>; снимаю сапоги. Через двадцать минут надо идти, непременно одному; выясняется, что дороги не знаю. Отбой. Новенький приехал еще где-то через час, без меня справились.

    \mds{2017-05-13, суббота}. Опять читаю за обедом и ужином. Гена, начальник наш, уехал на какое-то время, за него Сергей Голубев, чтец в обычное время. Кроме меня, теперь почти некому.

    Снова весь день кололи сорные дрова с гвоздями. Может, даже их не осталось больше.
    
    \mds{2017-05-14, воскресенье}. Рыбные котлеты на обед и на ужин, очень вкусные. Читаю утром, читаю вечером. Саша, вновь прибывший, ходит на все службы подряд, в келье редко увидишь. Спросил меня, где ноги можно помыть. Ответ: нигде. Стирка носков тоже нелегальна, все из-за ограниченной емкости отстойника.

    \mds{2017-05-15, понедельник}. До обеда~---~погрузка минеральной ваты в мешках на тракторный прицеп. После обеда~---~выгрузка этих мешков на чердак игуменского корпуса. Минеральная вата с виду ничем не отличается от стекловаты, но, как и обещали, далеко не такая зловредная. Иначе было бы очень сложно.

    \mds{2017-05-16, вторник}. Послушания заурядные: утром мытье полов второго этажа игуменского корпуса, после обеда разборка стены сарая и погрузка старой поленницы на прицеп. Незауряден конфликт по поводу времени окончания оказался. Виталий Желнов открыто поскандалил с Сергеем Голубевым на тему собирания бумажек на территории после сарая и дров,~---~чтоб не бездельничать до пяти вечера. Сказалось, думаю, и ползучее удлинение рабочего дня на полчаса~---~кто ж этого не заметит?

    Ну и дивная весточка из дому, продукт железной логики и глубокого проникновения в суть вещей. Запоминающийся день, в общем.

    \mds{2017-05-17, среда}. Рыба. Горбуша, крупная, одно удовольствие чистить. Моня лежит на травке в \emph{тени} от тачки. При плюс пяти жарко ему.

    \mds{2017-05-18, четверг}. До обеда разобрали крышу вторничного сарая, и еще собачью конуру рядом у стенки. Конура оказалась на редкость солидная, едва ли не крепче самого сарая. После обеда плавали на лодке через Пинегу, грузить сено на дальнем берегу. Вода в Пинеге темная, с торфяным оттенком, как в карельских озерах. Сено грузили в ту же самую лодку, сделавшую четыре рейса. Возвращались в монастырь через лес, видели муравейник с активными муравьями.

    Совсем поздно вечером видел в Казанском комара. Начинается.

\end{document}

